\section{Motivation}


An interactive or embodied system has to respond while the action is still
in progress. A closed label such as ``sitting'' is too coarse for that
response. Sitting still, sitting and looking at a phone, and sitting and
waving call for different responses, yet they collapse to one class if the
recognizer only has a small label set. We therefore describe the activity
with an English sentence. Running a strong vision-language model on every captured frame is impractical
for an interactive IoT system because the model is computationally expensive.
Our design therefore utilizes a small model on the smartphone which can write that sentence immediately, and a server with a stronger Vision-Language Model (VLM) can check the sentence and produce a more accurate description. However, sending every picture to the server makes the user wait on the uplink. The system when the picture has to be uploaded to the server. To achieve this, only a
low-confidence frame is JPEG-encoded and sent to the server, where the confidence of
that sentence is the minimum probability among the generated tokens. The server then
checks the picture and replaces the sentence when the two descriptions disagree.

Besides, checking a selected frame has an end-to-end confirmation deadline. That
deadline is shared by wireless transmission and server computation. A fixed
communication budget therefore cannot reliably guarantee an end-to-end
verification deadline. We therefore investigate a compute-aware communication scheduler. Before a
verification frame is uploaded, the system predicts its VLM latency
and subtracts the predicted computation time from the end-to-end deadline.
The remaining time becomes the communication budget of that frame. When
multiple clients request verification simultaneously, these dynamically
derived communication deadlines are used to determine their upload order.